\documentclass[11pt,a4paper]{article}
% Preâmbulo comum. Não duplicar pacotes nos módulos.
\usepackage[utf8]{inputenc}
\usepackage[T1]{fontenc}
\usepackage[brazil]{babel}
\usepackage{amsmath,amssymb,amsthm,mathtools}
\usepackage{geometry}
\geometry{a4paper,margin=2.3cm}
\usepackage{microtype}
\usepackage{enumitem}
\usepackage{graphicx}
\usepackage{xcolor}
\usepackage{booktabs}
\usepackage{hyperref}
\hypersetup{colorlinks=true,linkcolor=blue!40!black,urlcolor=blue!50!black}

\setcounter{tocdepth}{2}
\numberwithin{equation}{section}

\theoremstyle{plain}
\newtheorem{teorema}{Teorema}[section]
\newtheorem{proposicao}[teorema]{Proposição}
\newtheorem{lema}[teorema]{Lema}
\newtheorem{corolario}[teorema]{Corolário}

\theoremstyle{definition}
\newtheorem{definicao}[teorema]{Definição}
\newtheorem{exemplo}[teorema]{Exemplo}
\newtheorem{exercicio}{Exercício}[section]

\theoremstyle{remark}
\newtheorem{observacao}[teorema]{Observação}
\newtheorem{contraexemplo}[teorema]{Contraexemplo}
\newtheorem{intuicao}[teorema]{Intuição}

\newcommand{\R}{\mathbb{R}}
\newcommand{\N}{\mathbb{N}}
\newcommand{\Z}{\mathbb{Z}}
\newcommand{\Q}{\mathbb{Q}}
\DeclareRobustCommand{\eps}{\varepsilon}
\DeclareRobustCommand{\vet}[1]{\vec{#1}}
\DeclareRobustCommand{\dif}{\mathrm{d}}
\DeclareRobustCommand{\abs}[1]{\left\lvert #1 \right\rvert}
\DeclareRobustCommand{\norm}[1]{\left\lVert #1 \right\rVert}
\DeclareMathOperator{\grad}{grad}
\DeclareMathOperator{\Div}{div}
\DeclareMathOperator{\rot}{rot}

\newcommand{\base}[1]{\par\medskip\noindent\textbf{Base para relembrar.} #1\par\medskip}
\newcommand{\cuidado}[1]{\par\medskip\noindent\textbf{Quando a hipótese falha.} #1\par\medskip}

\newcommand{\figura}[2]{%
  \begin{figure}[htbp]
  \centering
  \IfFileExists{#1}{%
    \includegraphics[width=0.84\textwidth]{#1}%
  }{%
    \fbox{\parbox{0.8\textwidth}{\centering Figura ainda não gerada:\\ \texttt{#1}}}%
  }
  \caption{#2}
  \end{figure}%
}

\newenvironment{esqueleto}{%
  \par\medskip\noindent\textbf{Esqueleto da demonstração.}\par\nopagebreak
  \begin{enumerate}[leftmargin=1.4em]
}{\end{enumerate}\medskip}

\title{Gabarito: curvas paramétricas e coordenadas polares}
\author{}
\date{}
\begin{document}
\maketitle
\tableofcontents

Cada solução repete o enunciado e faz a conta até o fim. A numeração é a da nota \texttt{02\_parametricas\_polares.tex}.

\section{Básicos}

\begin{exercicio}[Básico --- B1]
Reta tangente de \(x=t^{2}\), \(y=2t\).
\end{exercicio}

\noindent\textbf{Solução.}
As derivadas são \(x'=2t\) e \(y'=2\), contínuas em \(\R\). Para \(t\neq 0\),
\[
\frac{\dif y}{\dif x}=\frac{y'}{x'}=\frac{2}{2t}=\frac{1}{t}.
\]
Em \(t=0\), \(x'=0\) e \(y'=2\neq 0\). A tangente é vertical, no ponto \(\bigl(x(0),y(0)\bigr)=(0,0)\). A reta é \(X=0\).

Em \(t=1\), o ponto é \((1,2)\) e a inclinação é \(1/1=1\). A reta tangente é
\[
Y-2=1\cdot(X-1),\qquad Y=X+1.
\]

\begin{exercicio}[Básico --- B2]
Comprimento de \(x=3\cos t\), \(y=3\sin t\), \(t\in[0,\pi/2]\).
\end{exercicio}

\noindent\textbf{Solução.}
Temos \(x'=-3\sin t\), \(y'=3\cos t\) e
\[
\sqrt{x'^{2}+y'^{2}}
=\sqrt{9\sin^{2}t+9\cos^{2}t}
=\sqrt{9}=3.
\]
O caminho é um quarto de circunferência de raio \(3\), percorrido uma vez. O comprimento é
\[
L=\int_{0}^{\pi/2} 3\,\dif t=3\cdot\frac{\pi}{2}=\frac{3\pi}{2}.
\]

\begin{exercicio}[Básico --- B3]
Área polar do disco de raio \(2\).
\end{exercicio}

\noindent\textbf{Solução.}
A região \(0\le\theta\le 2\pi\), \(0\le\rho\le 2\) é o disco de raio \(2\) centrado na origem. A função \(r(\theta)=2\) é contínua e não negativa, e o intervalo de ângulos tem comprimento \(2\pi\), varrendo o disco uma vez. A fórmula dá
\[
A=\frac12\int_{0}^{2\pi} 2^{2}\,\dif\theta
=\frac12\int_{0}^{2\pi} 4\,\dif\theta
=\frac12\cdot 4\cdot 2\pi
=4\pi.
\]
A área geométrica do disco é \(\pi R^{2}=\pi\cdot 4=4\pi\). Os dois valores coincidem.

\begin{exercicio}[Básico --- B4]
Revolução de \(x=t\), \(y=t\), \(t\in[0,2]\).
\end{exercicio}

\noindent\textbf{Solução.}
No intervalo, \(y=t\ge 0\). As derivadas são \(x'=1\) e \(y'=1\), e
\[
\sqrt{x'^{2}+y'^{2}}=\sqrt{1+1}=\sqrt{2}.
\]
A área da superfície é
\begin{align*}
S
&=\int_{0}^{2} 2\pi\cdot t\cdot\sqrt{2}\,\dif t
=2\pi\sqrt{2}\int_{0}^{2} t\,\dif t
=2\pi\sqrt{2}\Biggl[\frac{t^{2}}{2}\Biggr]_{0}^{2}
=2\pi\sqrt{2}\cdot\frac{4}{2}
=4\pi\sqrt{2}.
\end{align*}
O segmento de \((0,0)\) a \((2,2)\) gera um cone. A fórmula foi aplicada com \(y\ge 0\) e com derivadas contínuas.

\section{Intermediários}

\begin{exercicio}[Intermediário --- I1]
Comprimento de uma arcada da ciclóide.
\end{exercicio}

\noindent\textbf{Solução.}
Para \(x=t-\sin t\) e \(y=1-\cos t\),
\[
x'=1-\cos t,\qquad y'=\sin t.
\]
Então
\begin{align*}
x'^{2}+y'^{2}
&=(1-\cos t)^{2}+\sin^{2}t
=1-2\cos t+\cos^{2}t+\sin^{2}t
=2-2\cos t
=2(1-\cos t).
\end{align*}
A identidade \(1-\cos t=2\sin^{2}(t/2)\) dá
\[
2(1-\cos t)=4\sin^{2}(t/2),
\qquad
\sqrt{x'^{2}+y'^{2}}=2\abs{\sin(t/2)}.
\]
No intervalo \([0,2\pi]\), temos \(t/2\in[0,\pi]\) e \(\sin(t/2)\ge 0\). O módulo sai, e o integrando vale \(2\sin(t/2)\). Portanto
\begin{align*}
L
&=\int_{0}^{2\pi} 2\sin(t/2)\,\dif t
=2\cdot\bigl(-2\bigr)\Bigl[\cos(t/2)\Bigr]_{0}^{2\pi}
=-4\Bigl(\cos\pi-\cos 0\Bigr)
=-4\bigl(-1-1\bigr)
=8.
\end{align*}
A antiderivada de \(\sin(t/2)\) é \(-2\cos(t/2)\), e o fator \(2\) na frente produz o \(-4\). O comprimento de uma arcada é \(8\).

\begin{exercicio}[Intermediário --- I2]
Áreas da rosa \(r=\sin(2\theta)\).
\end{exercicio}

\noindent\textbf{Solução.}
Uma pétala em que \(r\ge 0\) é \(\theta\in[0,\pi/2]\): aí \(\sin(2\theta)\ge 0\), e \(r\) se anula nos extremos. A área é
\begin{align*}
A_{1}
&=\frac12\int_{0}^{\pi/2}\sin^{2}(2\theta)\,\dif\theta
=\frac12\int_{0}^{\pi/2}\frac{1-\cos(4\theta)}{2}\,\dif\theta
=\frac14\int_{0}^{\pi/2}\bigl(1-\cos(4\theta)\bigr)\dif\theta\\
&=\frac14\Biggl[\theta-\frac{\sin(4\theta)}{4}\Biggr]_{0}^{\pi/2}
=\frac14\Biggl(\frac{\pi}{2}-\frac{\sin(2\pi)}{4}-0\Biggr)
=\frac14\cdot\frac{\pi}{2}
=\frac{\pi}{8}.
\end{align*}
De \(0\) a \(2\pi\) a função \(\sin(2\theta)\) tem quatro trechos de comprimento \(\pi/2\) em que não é negativa, separados por zeros em \(0,\pi/2,\pi,3\pi/2,2\pi\). Integrar \(r^{2}\) não vê o sinal de \(r\):
\begin{align*}
A_{4}
&=\frac12\int_{0}^{2\pi}\sin^{2}(2\theta)\,\dif\theta
=\frac12\int_{0}^{2\pi}\frac{1-\cos(4\theta)}{2}\,\dif\theta
=\frac14\Biggl[\theta-\frac{\sin(4\theta)}{4}\Biggr]_{0}^{2\pi}
=\frac14\cdot 2\pi
=\frac{\pi}{2}.
\end{align*}
E \(4\cdot(\pi/8)=\pi/2\). As quatro pétalas têm área total \(\pi/2\), e cada uma tem área \(\pi/8\).

\begin{exercicio}[Intermediário --- I3]
Interseção de \(r=2\cos\theta\) e \(r=1\), e a área comum.
\end{exercicio}

\noindent\textbf{Solução.}
(i) Para \(r=2\cos\theta\ge 0\) restringimos \(\theta\in[-\pi/2,\pi/2]\). Igualando os raios,
\[
2\cos\theta=1\implies\cos\theta=\frac12\implies\theta=\frac{\pi}{3}\ \text{ou}\ \theta=-\frac{\pi}{3},
\]
ambos no intervalo. Os pontos cartesianos são
\begin{align*}
\theta=\frac{\pi}{3}:&\quad
\bigl(1\cdot\cos(\pi/3),\,1\cdot\sin(\pi/3)\bigr)
=\Bigl(\frac12,\frac{\sqrt{3}}{2}\Bigr),\\
\theta=-\frac{\pi}{3}:&\quad
\Bigl(\frac12,-\frac{\sqrt{3}}{2}\Bigr).
\end{align*}
O polo: \(2\cos\theta=0\) em \(\theta=\pm\pi/2\), então a curva \(r=2\cos\theta\) passa pela origem. A curva \(r=1\) tem raio constante \(1\) e não passa pela origem. O polo não é interseção.

(ii) A região dentro das duas curvas, em coordenadas polares com origem no centro do círculo \(r=1\). Para \(\theta\in[0,\pi/3]\), \(\cos\theta\ge\cos(\pi/3)=1/2\), logo \(2\cos\theta\ge 1\): ao longo do raio, o círculo \(r=1\) é a fronteira interna da interseção. Para \(\theta\in[\pi/3,\pi/2]\), \(2\cos\theta\le 1\): a fronteira interna é \(r=2\cos\theta\). O lado \(\theta\in[-\pi/2,0]\) repete a área por simetria. Assim
\begin{align*}
A
&=2\cdot\frac12\int_{0}^{\pi/3} 1^{2}\,\dif\theta
+2\cdot\frac12\int_{\pi/3}^{\pi/2}(2\cos\theta)^{2}\,\dif\theta
=\int_{0}^{\pi/3} 1\,\dif\theta
+\int_{\pi/3}^{\pi/2} 4\cos^{2}\theta\,\dif\theta.
\end{align*}
A primeira integral vale \(\pi/3\). Na segunda, \(\cos^{2}\theta=(1+\cos(2\theta))/2\), logo \(4\cos^{2}\theta=2\bigl(1+\cos(2\theta)\bigr)\), e
\begin{align*}
\int_{\pi/3}^{\pi/2} 4\cos^{2}\theta\,\dif\theta
&=\int_{\pi/3}^{\pi/2} 2\bigl(1+\cos(2\theta)\bigr)\dif\theta
=\Bigl[2\theta+\sin(2\theta)\Bigr]_{\pi/3}^{\pi/2}.
\end{align*}
No extremo superior: \(2\cdot(\pi/2)+\sin\pi=\pi\). No extremo inferior: \(2\cdot(\pi/3)+\sin(2\pi/3)=2\pi/3+\sqrt{3}/2\). A diferença é
\[
\pi-\frac{2\pi}{3}-\frac{\sqrt{3}}{2}
=\frac{\pi}{3}-\frac{\sqrt{3}}{2}.
\]
Somando a primeira integral,
\[
A=\frac{\pi}{3}+\frac{\pi}{3}-\frac{\sqrt{3}}{2}
=\frac{2\pi}{3}-\frac{\sqrt{3}}{2}.
\]
Os dois círculos têm raio \(1\): um centrado na origem, o outro, porque \(r=2\cos\theta\) equivale a \((x-1)^{2}+y^{2}=1\), centrado em \((1,0)\). A distância entre centros é \(1\). A área da interseção de dois discos de raio \(1\) com centros a distância \(1\) é este valor.

\section{Avançados}

\begin{exercicio}[Avançado --- A1]
Invariância do comprimento por reparametrização.
\end{exercicio}

\noindent\textbf{Solução.}
Seja \(t=\varphi(s)\), com \(\varphi:[c,d]\to[a,b]\) de classe \(C^1\), \(\varphi'>0\), \(\varphi(c)=a\) e \(\varphi(d)=b\). As funções compostas \(x(\varphi(s))\) e \(y(\varphi(s))\) são de classe \(C^1\). Pela regra da cadeia,
\[
\frac{\dif}{\dif s}x\bigl(\varphi(s)\bigr)=x'\bigl(\varphi(s)\bigr)\varphi'(s),
\qquad
\frac{\dif}{\dif s}y\bigl(\varphi(s)\bigr)=y'\bigl(\varphi(s)\bigr)\varphi'(s).
\]
O integrando do comprimento no parâmetro \(s\) é
\begin{align*}
&\sqrt{\bigl(x'(\varphi)\varphi'\bigr)^{2}+\bigl(y'(\varphi)\varphi'\bigr)^{2}}
=\sqrt{\bigl(\varphi'\bigr)^{2}\bigl(x'(\varphi)^{2}+y'(\varphi)^{2}\bigr)}
=\abs{\varphi'}\sqrt{x'(\varphi)^{2}+y'(\varphi)^{2}}.
\end{align*}
Como \(\varphi'>0\), o módulo é \(\varphi'\). Integramos:
\begin{align*}
\int_{c}^{d}\sqrt{x'(\varphi)^{2}+y'(\varphi)^{2}}\,\varphi'(s)\,\dif s.
\end{align*}
A substituição \(t=\varphi(s)\), \(\dif t=\varphi'(s)\,\dif s\), com extremos de \(a\) a \(b\), transforma a integral em
\[
\int_{a}^{b}\sqrt{x'(t)^{2}+y'(t)^{2}}\,\dif t.
\]
Os comprimentos coincidem. A hipótese \(\varphi'>0\) foi usada para tirar o módulo e para preservar a orientação dos extremos. Se \(\varphi'\) mudasse de sinal, a fórmula com \(\abs{\varphi'}\) ainda daria o comprimento do percurso, mas a substituição direta sem módulo inverteria pedaços.

\begin{exercicio}[Avançado --- A2]
A área polar de \(r=\cos\theta\) depende do intervalo.
\end{exercicio}

\noindent\textbf{Solução.}
Primeiro a geometria, para saber o que deveria dar a área. Com \(r=\cos\theta\),
\begin{align*}
x&=r\cos\theta=\cos^{2}\theta=\frac{1+\cos(2\theta)}{2},\\
y&=r\sin\theta=\cos\theta\sin\theta=\frac12\sin(2\theta).
\end{align*}
Então
\[
\Bigl(x-\frac12\Bigr)^{2}+y^{2}
=\frac14\cos^{2}(2\theta)+\frac14\sin^{2}(2\theta)
=\frac14.
\]
É a circunferência de centro \((1/2,0)\) e raio \(1/2\). A área do disco é \(\pi/4\).

(i) No intervalo \([-\pi/2,\pi/2]\), \(\cos\theta\ge 0\). Calculamos
\begin{align*}
\frac12\int_{-\pi/2}^{\pi/2}\cos^{2}\theta\,\dif\theta
&=\frac12\int_{-\pi/2}^{\pi/2}\frac{1+\cos(2\theta)}{2}\,\dif\theta
=\frac14\Biggl[\theta+\frac12\sin(2\theta)\Biggr]_{-\pi/2}^{\pi/2}.
\end{align*}
O seno de \(\pm\pi\) é zero, e
\[
\frac14\Biggl(\frac{\pi}{2}-\Bigl(-\frac{\pi}{2}\Bigr)\Biggr)
=\frac14\cdot\pi
=\frac{\pi}{4}.
\]
Este valor é a área do disco. O intervalo varre a circunferência uma vez, com raio não negativo.

(ii) De \(0\) a \(2\pi\),
\begin{align*}
\frac12\int_{0}^{2\pi}\cos^{2}\theta\,\dif\theta
&=\frac14\Biggl[\theta+\frac12\sin(2\theta)\Biggr]_{0}^{2\pi}
=\frac14\cdot 2\pi
=\frac{\pi}{2}.
\end{align*}
É o dobro de \(\pi/4\).

(iii) Para \(\theta\in(\pi/2,3\pi/2)\), \(\cos\theta<0\). O ponto de raio \(\cos\theta\) e ângulo \(\theta\) coincide com o ponto de raio \(-\cos\theta>0\) e ângulo \(\theta+\pi\). Explicitamente, o ponto de ângulo \(\theta+\pi\) e raio \(\cos(\theta+\pi)=-\cos\theta\) tem coordenadas
\begin{align*}
x&=(-\cos\theta)\cos(\theta+\pi)=(-\cos\theta)(-\cos\theta)=\cos^{2}\theta,\\
y&=(-\cos\theta)\sin(\theta+\pi)=(-\cos\theta)(-\sin\theta)=\cos\theta\sin\theta,
\end{align*}
as mesmas do ponto \((\cos\theta,\theta)\). Enquanto \(\theta\) percorre \((\pi/2,3\pi/2)\), a curva repete a circunferência já percorrida em \((-\pi/2,\pi/2)\). A fórmula \(\frac12\int r^{2}\) integra um quadrado, sempre não negativo, e conta o disco duas vezes. A hipótese ``cada ponto, salvo o polo, corresponde a um único \(\theta\)'' falha no intervalo \([0,2\pi]\).

\begin{exercicio}[Avançado --- A3]
Revolução com \(y\) negativo.
\end{exercicio}

\noindent\textbf{Solução.}
A integral pedida é
\begin{align*}
\int_{0}^{2\pi} 2\pi\sin t\,\dif t
&=2\pi\Bigl[-\cos t\Bigr]_{0}^{2\pi}
=2\pi\bigl(-\cos(2\pi)-(-\cos 0)\bigr)
=2\pi\bigl(-1-(-1)\bigr)
=0.
\end{align*}
Ela não é a área da esfera. Na fórmula
\[
S=\int_{a}^{b} 2\pi\, y(t)\,\sqrt{x'^{2}+y'^{2}}\,\dif t,
\]
aplicada a \(x=\cos t\), \(y=\sin t\), o fator \(\sqrt{x'^{2}+y'^{2}}\) vale \(1\), mas \(y=\sin t\) é negativo em \((\pi,2\pi)\). A hipótese \(y\ge 0\) falha. A integral trata a metade de baixo como raio negativo e subtrai a área que a metade de cima somou: \(2\pi\int_{0}^{\pi}\sin t\,\dif t=4\pi\) e \(2\pi\int_{\pi}^{2\pi}\sin t\,\dif t=-4\pi\). A área da esfera de raio \(1\) é \(4\pi\), obtida girando só o semicírculo superior, onde \(y\ge 0\).

\end{document}
